\documentclass[11pt,a4paper,openany]{memoir}
\usepackage{mathpazo}
\usepackage{fontspec}
\setmainfont{TeX Gyre Pagella}
\setsansfont{TeX Gyre Heros}
\setmonofont{TeX Gyre Cursor}
\usepackage{natbib}
\newcommand*{\olpath}{../..}
\newcommand*{\ollangid}{fr}
\newcommand*{\ollanguage}{french}
\input{\olpath/sty/open-logic.sty}
\input{\olpath/sty/open-logic-defer.sty}
\includeenv{editorial}
% Les cadres doivent restituer les notes après leur fermeture.
\usepackage{footnotehyper}
\makesavenoteenv{editorial}
\frcompletevariantstrue
% Les renvois hors livraison restent des liens vers leur section originale.
% Dès que la section française est intégrée, son renvoi interne est prioritaire.
\makeatletter
\newcommand{\FRoriginalsection}[2]{%
  section « \href{https://github.com/OpenLogicProject/OpenLogic/blob/9620cc73f9c8e0ad003c514a5d3748f29611c4c0/content/history/set-theory/#1.tex}{#2} »\space
  (dans l'original anglais, hors de cette livraison)}
\newcommand{\FRoriginalmethodsection}[2]{%
  section « \href{https://github.com/OpenLogicProject/OpenLogic/blob/9620cc73f9c8e0ad003c514a5d3748f29611c4c0/content/methods/induction/#1.tex}{#2} »\space
  (dans l'original anglais, hors de cette livraison)}
\csdef{FRexternal@his:set:mythology:sec}{%
  \FRoriginalsection{mythology}{More Myth than History?}}
\csdef{FRexternal@his:set:limits:sec}{%
  \FRoriginalsection{limits}{Rigorous Definition of Limits}}
\csdef{FRexternal@sth:::part}{%
  partie « \href{https://github.com/OpenLogicProject/OpenLogic/blob/9620cc73f9c8e0ad003c514a5d3748f29611c4c0/content/set-theory/set-theory.tex}{Set Theory} »\space
  (dans l'original anglais, hors de cette livraison)}
\csdef{FRexternal@sth:ord-arithmetic::chap}{%
  chapitre « \href{https://github.com/OpenLogicProject/OpenLogic/blob/9620cc73f9c8e0ad003c514a5d3748f29611c4c0/content/set-theory/ord-arithmetic/ord-arithmetic.tex}{Ordinal Arithmetic} »\space
  (dans l'original anglais, hors de cette livraison)}
\csdef{FRexternal@mth:ind:idf:sec}{%
  \FRoriginalmethodsection{inductive-definitions}{Inductive Definitions}}
\csdef{FRexternal@mth:ind:sti:sec}{%
  \FRoriginalmethodsection{structural-induction}{Structural Induction}}
\RenewDocumentCommand \olref { o o o m } {%
  \begingroup
  \edef\FRrefkey{%
    \IfNoValueTF{#1}
      {\theolpart:\theolchapter:\theolsection}
      {\IfNoValueTF{#2}
        {\theolpart:\theolchapter:#1}
        {\IfNoValueTF{#3}{\theolpart:#1:#2}{#1:#2:#3}}}:#4}%
  \@ifundefined{r@\FRrefkey}
    {\ifcsdef{FRexternal@\FRrefkey}
      {\csuse{FRexternal@\FRrefkey}}{\cref{\FRrefkey}}}
    {\cref{\FRrefkey}}%
  \endgroup}
\makeatother

\problemsperchapter
\setlrmarginsandblock{28mm}{28mm}{*}
\setulmarginsandblock{26mm}{28mm}{*}
\checkandfixthelayout
\linespread{1.05}
\setlength{\emergencystretch}{2em}
\let\cleardoublepage\clearpage
\pagestyle{plain}
\hypersetup{pdftitle={OpenLogic : édition française — Ensembles, fonctions et construction des nombres},
pdfauthor={Open Logic Project ; traduction française assistée par IA},
pdfsubject={Ensembles, relations, fonctions, dénombrabilité, nombres réels et ensembles infinis},
linkcolor=black,citecolor=black,urlcolor=blue}
\begin{document}
\frontmatter
\begin{titlingpage}
\setlength{\parindent}{0pt}
\vspace*{15mm}
{\sffamily\Huge\bfseries OpenLogic\par}
\vspace{6mm}
{\sffamily\LARGE Édition française\par}
\vspace{14mm}
{\sffamily\huge Ensembles, fonctions\\et construction des nombres\par}
\vspace{5mm}
{\large Relations, dénombrabilité et arithmétisation :\\
des premiers ensembles aux nombres réels et à l'infini\par}
\vfill
{\large Cinquième livraison : six chapitres complets\par}
\vspace{4mm}
Cette livraison comprend quarante-cinq sections et leurs six chapitres d’organisation,
soit 51 des 722 unités de contenu prévues pour l’édition intégrale.
Elle conserve les démonstrations, exemples, figures et exercices
des chapitres originaux. Les précisions éditoriales sont signalées
dans le texte. Les présentations alternatives de la dénombrabilité
sont réunies dans un même chapitre, avec leurs conventions explicites.
Le cinquième chapitre construit les entiers, les rationnels et les réels,
par les coupures de Dedekind et les classes de suites de Cauchy.
Le sixième étudie les ensembles infinis, les algèbres de Dedekind,
la récurrence et une autre preuve du théorème de Schröder--Bernstein.
\vspace{7mm}
\par\small
Texte original : \href{https://openlogicproject.org/}{Open Logic Project}.
Traduction française et révision éditoriale assistées par IA.
Aucune validation humaine indépendante n'est revendiquée.
\par\medskip
Texte original et traduction : licence
\href{https://creativecommons.org/licenses/by/4.0/}{Creative Commons
Attribution 4.0 International}. Les notices de droits propres aux
composants sont conservées avec les sources.
\par\medskip
\href{https://kokunoyumeto.github.io/OpenLogic-translations/}{Catalogue
des traductions d'OpenLogic}
\end{titlingpage}
\tableofcontents*
\mainmatter
\olimport[content/sets-functions-relations/sets]{sets}
\olimport[content/sets-functions-relations/relations]{relations-complete}
\olimport[content/sets-functions-relations/functions]{functions}
\olimport[content/sets-functions-relations/size-of-sets]{size-of-sets-complete}
\olimport[content/sets-functions-relations/arithmetization]{arithmetization}
\olimport[content/sets-functions-relations/infinite]{infinite}
\backmatter
% La notice Frege1884 cite elle-même cette traduction anglaise.
\nocite{Frege1953}
\bibliographystyle{plainnat}
\bibliography{../../bib/open-logic}
\end{document}
